\chapter{Hodge theory: the tools we use}\label{ch:toolkit}

This chapter collects what the rest of the book needs. Most of it is
standard and is proved in Voisin's books \cite{Voi02,Voi03} or in Griffiths
and Harris \cite{GH78}. We give proofs when they are short and when the
argument itself matters later.

\section{Hodge structures}

\begin{definition}
A \emph{rational Hodge structure of weight $k$} is a finite-dimensional
$\QQ$-vector space $V$ together with a decomposition
\[
  V_{\CC}=V\otimes_{\QQ}\CC=\bigoplus_{p+q=k}V^{p,q},\qquad
  \overline{V^{p,q}}=V^{q,p}.
\]
A \emph{morphism} of Hodge structures of the same weight is a
$\QQ$-linear map whose complexification preserves the decomposition.
\end{definition}

The Hodge filtration $F^{p}V_{\CC}=\bigoplus_{p'\ge p}V^{p',k-p'}$
determines the decomposition, through $V^{p,q}=F^{p}\cap\overline{F^{q}}$.
In families it is the filtration, not the decomposition, that varies
holomorphically.

Direct sums, tensor products, duals and exterior powers of Hodge
structures are again Hodge structures, with the obvious types; for example
$(V\otimes W)^{p,q}=\bigoplus V^{a,b}\otimes W^{p-a,q-b}$, and the dual of
a structure of weight $k$ has weight $-k$, with $(V^{\vee})^{p,q}$ dual to
$V^{-p,-q}$. The \emph{Tate structure} $\QQ(1)$ is the one-dimensional
structure of weight $-2$ and type $(-1,-1)$, and $V(m)=V\otimes\QQ(1)^{\otimes m}$
shifts the weight by $-2m$ without changing anything else.

\begin{definition}
A \emph{Hodge class} in a Hodge structure $V$ of even weight $2p$ is an
element of $V\cap V^{p,p}$. Equivalently, it is a morphism of Hodge
structures $\QQ(-p)\to V$.
\end{definition}

For a smooth projective $X$, the cohomology $H^{k}(X,\QQ)$ is a Hodge
structure of weight $k$ by the Hodge decomposition, and $\Hdg^{p}(X)$ is
the space of Hodge classes in $H^{2p}(X,\QQ)$. The second description is
the useful one: a Hodge class is the same thing as a morphism of Hodge
structures, and the Hodge conjecture predicts that such morphisms come
from geometry.

\section{Polarisations}

Let $V$ be a Hodge structure of weight $k$. The operator $C$ that acts on
$V^{p,q}$ by $i^{p-q}$ is defined over $\RR$, and it is called the Weil
operator.

\begin{definition}
A \emph{polarisation} of $V$ is a bilinear form $Q\colon V\times V\to\QQ$,
symmetric if $k$ is even and alternating if $k$ is odd, such that the
complexification satisfies $Q(V^{p,q},V^{p',q'})=0$ unless $p=q'$, and such
that the hermitian form $Q(Cx,\bar y)$ is positive definite on $V_{\CC}$.
\end{definition}

\begin{proposition}\label{prop:semisimple}
Let $V$ be a polarised Hodge structure and $U\subseteq V$ a sub-Hodge
structure, that is, a subspace with $U_{\CC}=\bigoplus(U_{\CC}\cap V^{p,q})$.
Then the orthogonal complement $U^{\perp}$ with respect to $Q$ is a sub-Hodge
structure, and $V=U\oplus U^{\perp}$. In particular, polarisable Hodge
structures form a semisimple category.
\end{proposition}

\begin{proof}
The form $h(x,y)=Q(Cx,\bar y)$ is positive definite, so $U_{\CC}$ and its
$h$-orthogonal complement span $V_{\CC}$ and meet in zero. Since $C$
preserves $U_{\CC}$, the $h$-orthogonal complement of $U_{\CC}$ equals the
$Q$-orthogonal complement of $\overline{U_{\CC}}=U_{\CC}$, which is
$(U^{\perp})_{\CC}$. So $V=U\oplus U^{\perp}$ over $\QQ$. The $Q$-orthogonal
of a subspace stable under the decomposition is stable under it, because
$Q$ pairs $V^{p,q}$ only with $V^{q,p}$.
\end{proof}

The cohomology of a smooth projective variety is polarisable. This is the
content of the Hodge--Riemann relations, recalled in
Section~\ref{sec:kahler}. Semisimplicity is used constantly: it lets us
split off the part of a cohomology group that we understand and look at
what is left.

\section{Mumford--Tate and Hodge groups}\label{sec:mt}

A Hodge structure of weight $k$ on $V$ is the same thing as a homomorphism
of real algebraic groups
\[
  h\colon\mathbb S=\operatorname{Res}_{\CC/\RR}\mathbb G_{m}\longrightarrow
  \GL(V_{\RR}),
\]
where $z\in\mathbb S(\RR)=\CC^{\times}$ acts on $V^{p,q}$ by
$z^{-p}\bar z^{-q}$. The weight is recorded by the restriction of $h$ to
$\RR^{\times}$, and the Weil operator is $h(i)$.

\begin{definition}
The \emph{Mumford--Tate group} $\MT(V)$ is the smallest algebraic subgroup
of $\GL(V)$, defined over $\QQ$, whose real points contain $h(\mathbb
S(\RR))$. The \emph{Hodge group} $\Hg(V)$ is the smallest one whose real
points contain $h(U(1))$, where $U(1)\subset\mathbb S(\RR)$ is the unit
circle.
\end{definition}

For a structure of weight $k\ne0$ the Mumford--Tate group is the product of
the Hodge group and the homotheties. The Hodge group is the right object
for counting Hodge classes.

\begin{theorem}\label{thm:hginvariants}
Let $V$ be a rational Hodge structure and let $T=V^{\otimes a}\otimes
(V^{\vee})^{\otimes b}$ be a tensor space, with its Hodge structure of some
weight $w$. If $w$ is even, then an element of $T$ is a Hodge class if and
only if it is invariant under $\Hg(V)$. If $V$ is polarisable, $\Hg(V)$ is
a connected reductive group.
\end{theorem}

\begin{proof}
On $T^{p,q}$ the circle acts through $z\mapsto z^{-p}\bar z^{-q}=z^{q-p}$,
since $\bar z=z^{-1}$ on $U(1)$. So the fixed vectors of $h(U(1))$ in
$T_{\CC}$ are exactly the elements of $\bigoplus_{p}T^{p,p}$, which is
$T^{w/2,w/2}$. An element $t\in T$ (rational) is fixed by $h(U(1))$ if and
only if it is fixed by the smallest $\QQ$-algebraic subgroup containing
$h(U(1))$, because the stabiliser of $t$ is an algebraic subgroup defined
over $\QQ$. This is the first claim. The group is connected because $U(1)$
is. Reductivity uses the polarisation: the form $Q(Cx,\bar y)$ makes the
real points of a certain inner form of $\Hg(V)$ compact, which forces
reductivity. We refer to \cite{Del82} for this step.
\end{proof}

So the Hodge ring of $X$, in so far as it sits inside tensor constructions
on one Hodge structure, is a ring of invariants. For abelian varieties all
of the cohomology sits inside the exterior algebra on $H^{1}$, and this
observation is the starting point of Chapter~\ref{ch:abelian}.

\section{The K\"ahler package}\label{sec:kahler}

Let $X$ be smooth projective of dimension $d$ and let $\eta\in H^{2}(X,\QQ)$
be the class of an ample divisor. Cup product with $\eta$ is the
\emph{Lefschetz operator} $L$. It is a morphism of Hodge structures
$H^{k}(X,\QQ)\to H^{k+2}(X,\QQ)(1)$, because $\eta$ is of type $(1,1)$.

\begin{theorem}[Hard Lefschetz]\label{thm:hardlefschetz}
For $0\le k\le d$ the map $L^{d-k}\colon H^{k}(X,\QQ)\to H^{2d-k}(X,\QQ)$ is
an isomorphism.
\end{theorem}

Define the primitive part $P^{k}=\ker\bigl(L^{d-k+1}\colon H^{k}\to
H^{2d-k+2}\bigr)$. Hard Lefschetz gives the \emph{Lefschetz decomposition}
\[
  H^{k}(X,\QQ)=\bigoplus_{j\ge0}L^{j}P^{k-2j},
\]
and every summand is a sub-Hodge structure, up to Tate twist.

\begin{theorem}[Hodge--Riemann bilinear relations]\label{thm:hodgeriemann}
For $k\le d$ the form
\[
  Q(x,y)=(-1)^{k(k-1)/2}\int_{X}\eta^{d-k}\cup x\cup y
\]
is a polarisation of the Hodge structure $P^{k}$. Twisting by the
appropriate signs on the summands of the Lefschetz decomposition gives a
polarisation of $H^{k}(X,\QQ)$.
\end{theorem}

Both theorems are proved with harmonic forms; see \cite[Chapter~6]{Voi02}.

\section{Correspondences}

By the K\"unneth formula and Poincar\'e duality,
\[
  H^{2d_{X}+2r}(X\times Y,\QQ)\cong\bigoplus_{k}\Hom\bigl(H^{k}(X,\QQ),
  H^{k+2r}(Y,\QQ)\bigr)
\]
for smooth projective $X,Y$ with $\dim X=d_{X}$. A class
$\Gamma\in H^{*}(X\times Y,\QQ)$ acts as
$\Gamma_{*}(x)=\mathrm{pr}_{Y*}(\Gamma\cup\mathrm{pr}_{X}^{*}x)$. Under this
isomorphism the Hodge classes on $X\times Y$ correspond to morphisms of
Hodge structures, up to Tate twist.

\begin{proposition}\label{prop:corr}
If $\Gamma$ is the class of an algebraic cycle on $X\times Y$, then
$\Gamma_{*}$ carries classes of algebraic cycles on $X$ to classes of
algebraic cycles on $Y$.
\end{proposition}

\begin{proof}
Pullback along the projection, intersection with a cycle and proper
pushforward are all defined on Chow groups, and they are compatible with
the cycle class map \cite[Chapter~19]{Ful98}.
\end{proof}

This is the basic mechanism for moving algebraicity around: once a single
correspondence is known to be algebraic, it transports algebraicity from
one variety to another. Most known cases of the Hodge conjecture for
abelian varieties are proved this way, starting from divisors.

\section{Degree \texorpdfstring{$2d-2$}{2d-2}}

\begin{corollary}\label{cor:degree2d-2}
If the Hodge conjecture holds for $X$ in degree $2$, it holds in degree
$2d-2$. In particular, by the theorem of Lefschetz in
Chapter~\ref{ch:lefschetz}, it holds in degree $2d-2$ for every smooth
projective $X$.
\end{corollary}

\begin{proof}
The map $L^{d-2}\colon H^{2}(X,\QQ)\to H^{2d-2}(X,\QQ)(d-2)$ is an
isomorphism of Hodge structures by Theorem~\ref{thm:hardlefschetz}. An
isomorphism of Hodge structures is a bijection on Hodge classes. So every
Hodge class of degree $2d-2$ has the form $\eta^{d-2}\cup\alpha$ with
$\alpha\in\Hdg^{1}(X)$. If $\alpha$ is the class of a divisor $D$, then
$\eta^{d-2}\cup\alpha$ is the class of the curve $D\cdot H_{1}\cdots
H_{d-2}$ for general hyperplane sections $H_{j}$ in the class of $\eta$.
\end{proof}

The argument is worth looking at closely because it shows the general
pattern. The Lefschetz operator $L$ is algebraic, being intersection with
a hyperplane. Its inverse on the image is not known to be algebraic. That
is the Lefschetz standard conjecture, and if it were known, the Hodge
conjecture in degree $2p$ for $p>d/2$ would follow from degree
$2d-2p$. The problem is to find inverses, not to find the forward map.
This pattern will appear again in Chapter~\ref{ch:remaining}.

\section{Families}\label{sec:families}

Let $\pi\colon\mathcal X\to S$ be a smooth projective morphism with connected
base. The groups $H^{k}(X_{s},\QQ)$ form a local system $R^{k}\pi_{*}\QQ$, and
the Hodge filtrations vary holomorphically and satisfy Griffiths
transversality. The result is a variation of Hodge structure. A Hodge class
at one point need not remain a Hodge class when it is transported to
nearby points. The set of nearby points where it does is called the
\emph{Hodge locus} of the class.

\begin{theorem}[Cattani--Deligne--Kaplan \cite{CDK95}]\label{thm:cdk}
For a variation of Hodge structure coming from a smooth projective family
over a quasi-projective base, the Hodge locus of every class is a countable
union of closed algebraic subvarieties of $S$. The locus of classes of
bounded norm (with respect to the polarisation) is a finite union.
\end{theorem}

The Hodge conjecture predicts this: if every Hodge class is algebraic, then
the Hodge locus is the image of a countable union of relative Hilbert
schemes, which are algebraic. Cattani, Deligne and Kaplan proved it with no
conjecture at all. Their theorem is one of the best pieces of evidence for
the Hodge conjecture.

\begin{theorem}[Deligne, global invariant cycles \cite{Del71}]\label{thm:invcycles}
Let $\bar{\mathcal X}$ be a smooth projective compactification of
$\mathcal X$. Then the classes in $H^{k}(X_{s},\QQ)$ invariant under the
monodromy of $\pi$ are exactly the restrictions of classes in
$H^{k}(\bar{\mathcal X},\QQ)$. Moreover, the invariant classes form a
sub-Hodge structure.
\end{theorem}

This theorem gives the standard reduction of the variational question to
the Hodge conjecture for the total space. Suppose a monodromy-invariant
class is a Hodge class at one point. By the second part of the theorem it
is then a Hodge class at every point. By the first part it is the
restriction of a class on $\bar{\mathcal X}$, and since restriction is a
surjective morphism of polarisable Hodge structures onto the invariant
part, Proposition~\ref{prop:semisimple} lets us choose that class to be a
Hodge class. If the Hodge conjecture
holds for $\bar{\mathcal X}$, the class is algebraic at every point. The
difficulty is that $\bar{\mathcal X}$ is a larger variety, about which we
know less.
